\subsection*{Theological and Spiritual Synthesis}

The peaceful king promised to Zion and the meek Son who reveals the Father disclose a form of divine rule in which power serves peace and authority teaches the burdened. Psalm 145 answers with praise of the gracious sovereign, so meekness appears as the exercise of faithful strength rather than its absence.

Romans gives this discipleship an interior and bodily depth. The Spirit who raised Jesus dwells in believers and enables them to put sinful deeds to death without treating the body itself as evil. Receptive littleness, Spirit-enabled obedience, and Christ's fitting yoke belong together as graced participation rather than self-salvation.
The celebration looks toward universal peace, resurrection, saving reception, and enduring praise. Communion may voice either Psalm 34's invitation to taste and seek refuge or Christ's summons to the burdened; both paths direct the assembly toward life received from God and lived under the gentle kingship of the Son.

\begin{fourstable}{Dimension}{Appointed anchors}{Source-grounded synthesis}{Limit}
Christological & Zech.~9:9--10; Ps.~145; Matt.~11:25--30 & Jesus discloses the peaceful King's power as filial revelation, meekness, and a yoke fitted by the Son who knows the Father. & The Lectionary officially correlates Zechariah and Matthew, and Zion remains the oracle's own horizon.\\
Pneumatological and anthropological & Rom.~8:9, 11--13 & The indwelling Spirit joins belonging to Christ, mortification of sinful deeds, and hope for the life of mortal bodies. & Romans remains semi-continuous; flesh is not bodily substance, and the proclamation omits verse 10.\\
Ecclesial and sacramental & Entrance; Collect; Prayer over Offerings; Communion A/B; Prayer after Communion & A people gathered by mercy is raised through the Son's humility, purified in offering, fed through a real alternative Communion path, and ordered toward salvation and praise. & The Week XIV orations serve A/B/C, and the Communion choice is unresolved.\\
Moral and eschatological & Gospel; Second Reading; Psalm; final prayer & Receptive littleness, meek strength, and Spirit-enabled obedience tend toward universal peace, bodily resurrection, salvation, and enduring praise. & Christ's yoke does not commend passivity before abuse, and promised rest is not present exemption from duty or suffering.\\
\end{fourstable}
